\documentclass[5pt, a4paper]{article}

\usepackage[utf8]{inputenc}
\usepackage[T1]{fontenc}
\usepackage{textcomp}
\usepackage{amsmath, amssymb}
\usepackage{multicol}
\usepackage{proof}
\usepackage{enumitem}
\usepackage{fancyhdr}
\usepackage[margin=1in]{geometry}
% figure support
\usepackage{import}
\usepackage{xifthen}
\pdfminorversion=7

\usepackage[most]{tcolorbox}
\usepackage{pdfpages}
\usepackage{transparent}
\newcommand{\incfig}[1]{%
    \def\svgwidth{\columnwidth}
    \import{./figures/}{#1.pdf_tex}
}

\tcbset{
colback=gray!5,
colframe=blue!50!white,
coltitle=black,
arc=0pt,
boxrule=-0.1mm,
fonttitle=\bfseries,
coltitle=black,
left=2.5mm,
leftrule=1mm,
right=3.5mm,
colbacktitle=black!10!white,
toptitle=0.75mm,
bottomtitle=0.5mm,
}
\begin{document}
\section{Eigenvalue Problems}
\subsection*{Recap}
    To recap, for a square matrix $A$ and a non-zero vector $x$: 
    \begin{align*}
        Ax &= \lambda x \\
        0 &= (A - \lambda I)x\\
        p(\lambda) &= \det(A - \lambda I) = 0
    .\end{align*}
    The solution to  $x$ is called the Eigenvectors, and solution to $\lambda$
    are the Eigenvalues. $p(\lambda)$ is called the characteristic polynomial.
    $A - \lambda I$ is singular. This leads to the spectral theorem for symetric
    real matrices, which states that 
    \begin{tcolorbox}[title=Spectral Theorem for Real Matrices]
        A Symetric real matrix $A$ can be diagonalized using an orthonormal
        basis of vectors. \[
        A = Q D Q^T
        .\] 
        $A$ is said to be orthogonally diagonalizable.
    \end{tcolorbox}
    This means that for any symetric real matrix, there always exists an
    orthonormal basis ($Q$) such that A is a simple diagonal matrix.
    Intuitively, this is changing the basis for $A$, such that $A$ becomes
    simple scaling.

    \subsection*{Polynomials and Eigenvalues}

    In numerical computation, the bulk of finding the root of a polynomial can
    be computed by computing the Eigenvalues, and such problem is classified
    as an eigenvalue problem. We do this by computing the eigenvalues using teh
    companion matrix, which is just the matrix form of the (monic, ie. the
    coeffecient for the largest term is $1$) characteristic polynomial, given by
    \[
    C_n = 
    \left(
    \begin{matrix}
        0 & 0 & \cdots &0& -c_0 \\
        1 & 0 & \cdots &0& -c_1 \\
        0 & 1 & \cdots &0& -c_2 \\
        \vdots & \vdots & \ddots & \vdots & \vdots \\
        0 & 0 & \cdots &1& -c_{n-1} \\
    \end{matrix}
    \right)
    .\] 

    For small polynomial ($n \le 4$), we can just use the readily available
    formula to compute the root. The Abel-Ruffini Theorem states that there is
    no general formula for finding the root of polynomial of power $\ge  5$
    (TODO: Read), which is very common, and we would need to rely on algorithm
    to solve it iteratively, and thus find the root of the polynomial
    indirectly. Generally, such algorithm solves the root one at a time, so
    calculating all the root at once usually boils down to solving the
    eigenvalues. 

    Before doing so, let's take a look at the question of conditioning, uniqueness and
    existence of eigenvectors and eigenvalues in the real field. In general, 
    \begin{enumerate}
        \item For an upper triangular matrix, the diagonal entries is the
            eigenvalues.
        \item All symetric matrices have atleast one eigenvector and
            eigenvalues.
        \item Every scalar multiple of an eigenvector is an eigenvector.
        \item Every eigenvalues of a matrix is unique.
         \item The algebraic multiplicity of a matrix $A$ is $det(A - \lambda I)$.
        \item The geometric multiplicity $dim(A - \lambda I)$ is equal to the
            number of unique eigenvalues of  $A$. It is equal or less than the
            algebraic multiplicity. If it is strictly less, the matrix $A$ is
            defective. Non-defective matrices of size $n\times n$ (ie. those with
            $n$ linearly independent eigenvectors) are diagonalizable. \[
                AX = XD \iff X^{-1} A X = D \iff A = XDX^{-1}
            ,\] 
            where $X$ is a matrix of eigenvectors, and $D$ is a diagonal matrix
            of the eigenvalues. If all the collumns in $X$ is orthogonal, then
            we say that $A$ is orthonormaly diagonalizable, and by the spectral
            theorem, $A = A^T$. More often than not, such matrices are symetric.
    \end{enumerate}

    Operations on a matrix changes the eigenvalues. In general, the following
    applies: 
    \begin{enumerate}
        \item Shifts $Ax = \lambda x \to (A- \sigma I)x = (\lambda -\sigma)x $
        \item Inverses $Ax = \lambda x \to  A^{-1}x = \frac{1}{\lambda}x$
        \item Powers $Ax = \lambda x \to A^k x = \lambda^k x$
        \item Polynomials $Ax = \lambda x \to P(A)x = P(\lambda)x$
        \item Cholseky, QR, Elemination does nto preserve the non-zero
            eigenvalues.
    \end{enumerate}
    The condition number (2-norm) for a non-defective eigenvalue problem is
    $X^{-1} A X = D$ is given by the condition numebr (2-norm) of $X$.
    \subsection*{Singular Value Decomposition}
    \begin{tcolorbox}[title=Singular Value Decomposition]
        For $A \in \mathbb{R}^{m \times n}$, \[
            A = U \Sigma V^T
        ,\] 
        where $U \in \mathbb{R}^{m \times m}$ and $V \in \mathbb{R}^{n \times
        n}$ are orthogonal. The matrices $A^{T}A$ and $AA^{T}$ are real symetric
        and orthogonally diagonalizable. $\Sigma$ is a diagonal matrix consisting
        of the singular value of A (the square root of $AA^{T}$ and $A^{T}A$
        non-zero eigenvalues.)
    \end{tcolorbox}
    When decomposing a matrix with SVD, we typically try to solve either
    $A^{T}A$ or $AA^{T}$ for their eigenvalues and eigenvectors. Typically, we
    choose whichever is least complex to compute.

    After calculating the eigenvalues and eigenvector, and forming either $U$ or
     $V$, we would need to form the other base. To see this, recall that
     $A^{T}A$ is positive semi-definite, and can be written as \[
         A^{T}A = VDV^{T} = V(\Sigma^{T} \Sigma)V^{T}.
    \]
    For each column of $V$ we have  $A^{T}A v_i = \sigma_i^2 v_i$. Multiplying
    both side with $A$ yields 

     \subsection*{Rootfinding Algorithm}
     \subsection*{Power Iteration and Variants}
     The idea behind power iteration is quite simple: For a matrix $A$ that we
     would like to compute the eigenvalues and eigenvectors, we start with a
     vector $v_0$ (random, most like,) and repeatedly applying $Av_i$
     iteratively. 

     If  $A$ is non-defective and the largest eigenvalue is strictly larger than
     the other ones, we can analyze the algorithm's convergence quite easily: 
\end{document}
